% ============================================================
% Page 5 — Alban aux « siècles bénédictins »
% Transcription validée du document source.
% Mise en page simple ; typographie et lettrines à traiter
% ultérieurement dans le préambule et le modèle général.
% ============================================================

\thispagestyle{plain}

\vspace*{12mm}

\begin{center}
    {\Large\bfseries Alban aux « siècles bénédictins »}
\end{center}

\vspace{18mm}

\noindent
Au VIIème siècle se fonde à Castres le monastère de Saint Benoît.
Un jour, notre petite cité sera, au spirituel, prise en charge par
ces Bénédictins. Alban sera ce qu'on appelle un « prieuré-cure ».
Le titulaire en sera le prieur. Le prieur en titre recevra les revenus
de la paroisse où un jour, il se fera remplacer par un « vicaire
perpétuel » dont il assurera l'entretien par un traitement appelé
« congrue ». En 1545, le monastère disparaîtra et c'est le chapitre
cathédral de Saint Benoît qui nommera le prieur d'Alban et assurera
sa subsistance.

\par\vspace{12mm}

\noindent
Un rarissime parchemin de 1231, un minuscule écrit de deux lignes
garantissant l'authenticité d'une relique -- une pierre sacrée --
sont les plus anciens vestiges de cette époque bénédictine. Des croix
de chemin (le calvaire de Notre Dame d'Outiguet est de 1750), les
croix de consécration de l'église (N.D. d'Alban), les armoiries de
la baronnie d'Alban, rappellent la croix bénédictine (qui s'inscrit
dans un cercle). Il est bien possible que St Michel de Nougaret,
St Michel de la Besse, aujourd'hui disparus aient semblable origine,
mais les documents font défaut.

\vfill
